\section{数据与实验设置}
\label{sec:data}

\paragraph{数据来源与规模}
空间数据为GSE251950，10x Visium，共10张切片\cite{gut2025}。GC1--GC9为原发灶，质控后点位数为1864--4256。GC6-PM（GSM7990479）在记录中写为胃癌转移灶，部位没有写成腹膜，因此单独列表，不参与9例排序。单细胞数据为GSE183904的处理后矩阵\cite{kumar2022}。其中腹膜肿瘤3份（sample19、sample20、sample38），原发肿瘤26份。sample36与sample38来自同一患者，分别是原发灶和腹膜灶。

\paragraph{划分}
排序只用9例原发灶。GC6-PM用同一套半径和评分定义计算，结果不计入5/9。腹膜比较用3份对26份的样本均值中位数，以及上述一对配对样本的$\log_2$倍数。置换在每张切片内部进行，切片之间不合并点位。

\paragraph{预处理}
保留组织内点位。去掉基因数低于200或线粒体比例高于25\%的点位。GC1由3834个点位减到3212个，其余切片几乎没有因这一阈值损失点位。计数归一化到每个点位1万，再取$\log(1+x)$。成纤维评分基因为DCN、LUM、COL1A1、FAP，巨噬评分基因为CD68、CSF1R、C1QA、C1QB。高分区是各切片内评分的前25\%。这8个基因不进入配体受体乘积。
